\chapter{RDMA at a Low Level}
\label{ch:lowlevel}

The history in \cref{sec:history} ended with a programming model that
outlived the fabric it was designed for, and the research questions of the
previous chapter depend on what that model allows. It is explained here from
the bottom up, because every design decision in the next chapter depends on what the network card does with a request and in which order. We
begin with the cost of an ordinary socket and arrive at remote direct memory
access and the verbs interface (\cref{ch:rdma}). We then introduce
libfabric, the portable interface on which the middleware is written,
together with its origins and the organizations behind it
(\cref{ch:libfabric}), follow a single transfer step by step in both
interfaces (\cref{ch:process}), and close with Ultra Ethernet, the transport
that has adopted the same interface (\cref{sec:uec}).

\section{From Sockets to RDMA and Verbs}
\label{ch:rdma}

The middleware moves data by remote direct memory access. The path to RDMA
starts at the ordinary socket, passes through the objects that a program
must create, and ends at three facts about ordering on which every later
design decision depends.

\subsection{The Cost of a Socket}

A program that sends data through a TCP socket hands a buffer to the kernel
with a system call. The kernel copies the bytes into its own buffers, cuts
them into segments, adds headers, and gives them to the network card's
driver. On the other machine an interrupt wakes the kernel, which checks the
segments, reorders them, copies the bytes into the receiving program's buffer
and wakes the program. Every message pays for two transitions into the
kernel, at least two copies, and the scheduler on the receiving side. For a
4 MiB block these costs are small next to the wire time. For an 8-byte credit
they are nearly all of the cost.

\subsection{Remote Direct Memory Access}

RDMA removes those costs by changing who does the work. The application
registers a region of its own memory with the network card and receives a key
for it. A peer that holds the key and the address can ask its own card to
write into that region, or read from it. The transfer is carried out by the
two cards. The kernel is not entered on the data path, the data is not copied
on the way, and for these one-sided operations the receiving application does
not run at all \citep{snia2025rdma,mellanox_ibwp}.

The idea is older than the network on which it became standard: the SHMEM
routines of the CRAY T3D already let one processor read or write the memory
of another without the remote processor taking part
\citep{feind1995shmem}, and InfiniBand, which was developed with the Virtual
Interface Architecture in mind, made that operation a service of the network
adapter \citep{mellanox_ibwp}. RoCE carries the
InfiniBand transport over Ethernet, and its second version, RoCEv2, wraps each
packet in UDP so that it can be routed like any other traffic
\citep{zhu2015dcqcn}. Soft-RoCE implements the RoCEv2 transport in software,
in the kernel, on top of any Ethernet card \citep{ibta_softroce2015}. It keeps
the programming interface and the semantics, and gives up the hardware
offload: the kernel builds and parses the packets after all. The lab pair runs on Soft-RoCE, and the home pair on cards that implement
RoCEv2 in hardware (\cref{ch:servers}).

\subsection{Objects}

The programming interface to RDMA on Linux is the verbs interface, provided
by rdma-core \citep{rdmacore}. A program creates five kinds of object before
a byte moves. The \emph{device context} is an open handle on one RDMA device, such as
\code{rxe0}. The \emph{protection domain} is a container that ties memory
regions and queue pairs together, so that a key from one domain is useless in
another. A \emph{memory region} is a registered range of the program's
memory: registering pins the pages and returns a local key and a remote key,
and a peer needs the remote key and the address to reach the region. A
\emph{queue pair} is the endpoint of a connection and consists of a send
queue and a receive queue. Work requests posted to the send queue are carried
out in order, and a reliable connected queue pair, the only kind used here,
retransmits lost packets and delivers in order. Finally, the
\emph{completion queue} is where the card reports that a work request
finished, with a status, and a program learns the outcome of every operation
by reading this queue.

A queue pair is created in a reset state and has to be walked through three
transitions before it can send: to an initialized state, in which its port and
access rights are set; to ready-to-receive, in which it learns the peer's
queue pair number, address and starting sequence number; and to
ready-to-send. The two sides need each other's numbers to make these
transitions, and RDMA has no way to exchange them before the connection
exists, so they travel over an ordinary TCP connection first. In this
project that exchange is a fixed 64-byte record
(\cref{sec:reused}). \Cref{fig:qpstates} shows the states.

\begin{figure}[htbp]
  \centering
  \begin{tikzpicture}[diagram,
      alt={State diagram of a queue pair. Four states in a row joined by arrows left to right:
        reset, initialized, ready to receive, ready to send. The arrow from reset to initialized
        is labeled set port and access rights. The arrow from initialized to ready to receive
        is labeled learn the peer's queue pair number, address and sequence number. The arrow
        from ready to receive to ready to send is labeled set timeouts and retry counts. A
        fifth state, error, sits below; dashed arrows lead to it from ready to receive and ready
        to send, labeled on a failed operation, and a dashed arrow leads from error back to
        reset.},
      st/.style={tint=#1, minimum width=25mm, minimum height=11mm, font=\footnotesize, align=center}]
    \node[st=dGray]   (r)   at (0,0)    {Reset};
    \node[st=dBlue]   (i)   at (4.4,0)  {Initialized};
    \node[st=dBlue]   (rtr) at (8.8,0)  {Ready to\\receive};
    \node[st=dGreen]  (rts) at (13.2,0) {Ready to\\send};
    \node[st=dOrange] (e)   at (11.0,-2.8) {Error};
    \draw[flow] (r) -- (i) node[midway, above=3pt, small text, align=center] {port,\\access};
    \draw[flow] (i) -- (rtr) node[midway, above=3pt, small text, align=center] {peer's\\numbers};
    \draw[flow] (rtr) -- (rts) node[midway, above=3pt, small text, align=center] {timeouts,\\retries};
    \draw[reply] (rtr.south) -- (e.north west);
    \draw[reply] (rts.south) -- (e.north east) node[midway, right=4pt, small text] {failed operation};
    \draw[reply] (e.west) -| (r.south);
  \end{tikzpicture}
  \caption[States of a queue pair]{States of a queue pair. Solid arrows are the path to a working connection; dashed arrows are failure and recovery.}
  \label{fig:qpstates}
\end{figure}

\subsection{Operations}

Four operations matter here. \emph{Send and receive} form the two-sided
pair: the receiver must have posted a receive buffer in advance, the send
consumes it, and both sides receive a completion. An \emph{RDMA write} is
one-sided. The sender names a remote address and key, and the data lands
there, while the receiving program receives no completion and is not informed
that anything arrived. An \emph{RDMA write with immediate data} is a write
that also carries 32 bits chosen by the sender. It consumes a posted receive
on the target, so the target receives a completion that carries those 32
bits, and it is the operation on which the blob channel is built: the data
goes directly to its slot, and the 32 bits identify the slot. An
\emph{RDMA read} is one-sided in the other direction, as the initiator pulls
bytes from a remote region.

RDMA write and read exist only on connected, reliable queue pairs. The
unreliable datagram service has neither. This is why both traffic classes in
the design use reliable connections (\cref{ch:compat}).

\subsection{Three Facts About Ordering}
\label{sec:ordering}

For a reliable connected queue pair, three facts about ordering follow from the specification, and the following chapters refer to them by number.

\begin{enumerate}
  \item Operations execute in order. Work requests posted to one send
        queue are carried out at the target in the order they were posted.
  \item Placement within one message is not ordered. A large write is
        split into packets, and nothing promises that the bytes of the message
        are placed in memory front to back. A reader that looks at the target
        memory while a write is in progress can see the end of the message
        before the middle.
  \item A completion follows placement. When the target receives the
        completion of a write with immediate data, all bytes of that write are
        in place.
\end{enumerate}

Fact 2 is the one most easily overlooked, as it means that the receiver must
not decide that a block has arrived by inspecting the block. The receiver
needs a signal that comes after the data, and facts 1 and 3 provide two ways
of obtaining one, of which \cref{sec:signals} compares five variants.

\subsection{Failure Modes, with a Production Reference}

RDMA moves the difficult parts of networking into the application's hands: memory
lifetimes, keys, connection state and completion handling. A parallel file
system is a useful reference for the consequences at scale. IBM Storage Scale,
the file system that began as GPFS \citep{schmuck2002gpfs}, uses RDMA on its
data path \citep{wuerthner2021gpfsrdma}, and its published issues include
a race between the RDMA layer and the file system that returned incorrect
data to a client \citep{ibm_ij44629} and a fault in reconnecting after a
failure \citep{ibm_iv61625}. Three rules follow from these reports, and the middleware applies each of
them: every completion is checked, a connection that enters the error state
is treated as lost and reported, and the arrival of data is never inferred
from anything except a signal that the transport orders after the data.

\section{libfabric}
\label{ch:libfabric}

The verbs interface serves RDMA devices and nothing else. Code written
against it cannot run on a laptop, cannot fall back to TCP, and cannot be
compared with a transport used by the connector of the serving engine itself.
The middleware is therefore written against libfabric
\citep{grun2015ofi,libfabric_repo}, the library of the OpenFabrics
Interfaces, and retains a second implementation on raw verbs so that the cost
of the portable layer can be measured. Its origins and the organizations
that contribute to it come first, then the reasons for adopting it, and
finally a map of its objects onto those of verbs.

\subsection{Origins of libfabric}
\label{sec:ofi-history}

The verbs interface predates the problem addressed here. It came
out of the InfiniBand specification, and the software stack of the
OpenFabrics Alliance began as an implementation of it, then grew to cover
iWARP and RoCE as those appeared \citep{grun2015ofi}. By the early 2010s the
users of that stack were no longer only people programming InfiniBand cards.
Libraries for message passing, for partitioned global address spaces, for
databases and for non-volatile memory all sat on top of verbs, and each had
to translate what it wanted into queue pairs and work requests.

A meeting at the Supercomputing conference in 2013 led to the OpenFabrics
Interfaces Working Group. The group asked those communities what they needed
from a network interface and collected 168 requirements. The design that came
out of the exercise divides the world in two: applications, which state what they need, and providers, which state what they can do, with a discovery step
that matches one to the other \citep{grun2015ofi}. The library that carries that interface is libfabric. It was designed under the alliance by companies, universities and national
laboratories together, and it remains a project of the alliance, maintained
by the working group \citep{grun2015ofi,xiong2025ofaue}.

\Cref{tab:ofi-timeline} places the versions relevant to this work in that
history, with release dates from the release notes in the repository of the
library \citep{libfabric_repo}.

\begin{table}[htbp]
  \centering
  \caption[libfabric from its start to the versions used here]{libfabric from its start to the versions used here.}
  \label{tab:ofi-timeline}
  \small
  \begin{tabular}{L{3.3cm}L{10.9cm}}
    \toprule
    Date & Event \\
    \midrule
    2013 & a meeting at the Supercomputing conference leads to the working group \\
    May 3, 2015 & libfabric 1.0.0 \\
    November 22, 2019 & libfabric 1.9.0, the oldest interface version the pool accepts \\
    December 16, 2022 & libfabric 1.17.0, the version on the development machine \\
    December 13, 2024 & libfabric 2.0.0, the version the Ultra Ethernet specification maps \\
    June 11, 2025 & Ultra Ethernet Specification 1.0 names libfabric as its interface (\cref{sec:uec}) \\
    September 18, 2026 & libfabric 2.7.0, the latest release when the repository was read, and the second version on which the tests of the pool run \\
    \bottomrule
  \end{tabular}
\end{table}

\subsection{Providers and the Organizations Behind Them}
\label{sec:ofi-providers}

The clearest record of who takes part in libfabric is the list of providers
in its source tree, because most providers are written by the maker or the
operator of the network they serve. \Cref{tab:providers} lists the providers
that the repository's description file presents, in the terms it uses
\citep{libfabric_repo}. Two further directories of the provider tree,
\code{coll} and \code{util}, are not described there and are left out.

\begin{table}[htbp]
  \centering
  \caption[Providers in the libfabric source tree]{Providers described in the libfabric source tree, and their role in this work.}
  \label{tab:providers}
  \small
  \begin{tabular}{L{2.9cm}L{6.6cm}L{4.7cm}}
    \toprule
    Provider & Network served & Role in this work \\
    \midrule
    verbs & InfiniBand, iWARP and RoCE devices, through the Linux verbs interface & the provider for both pairs on the cluster \\
    tcp, udp & any machine with sockets & tcp for development and tests \\
    sockets & a general provider over sockets, now deprecated in favor of tcp, udp and the utility providers & not used \\
    shm, sm2 & shared memory on one machine & shm for tests and for E7 \\
    rxm, rxd & utility providers that offer reliable-datagram endpoints over connected or datagram ones & rxm in E4 \\
    efa & Amazon EC2 Elastic Fabric Adapter & not used \\
    cxi & the Cray Slingshot network & not used here; the first choice of NIXL \\
    opx, psm2 & Omni-Path fabrics & not used \\
    psm3 & Ethernet and verbs datagram devices, with extra features on Intel E810 cards & not used \\
    usnic & Cisco VIC hardware in Cisco UCS servers & not used \\
    ucx & the UCX library, on NVIDIA InfiniBand fabrics & not used \\
    lpp & FabreX PCIe networks & not used \\
    mrail, lnx, hook & utility providers for several rails, for linking providers, and for intercepting calls & not used \\
    \bottomrule
  \end{tabular}
\end{table}

Three rows bear on the comparisons made in later chapters. UCX is both a
library in its own right, the one NIXL uses by default, and a provider
underneath libfabric. The cxi provider is the native interface of Slingshot, an Ethernet-compliant
network built from the Rosetta switch and the Cassini network card
\citep{libfabric_repo}, so a program written against libfabric runs on the
network of the largest supercomputers with no verbs layer in between.
Conversely, a library written for those networks may know nothing about the
verbs provider that commodity RoCE cards need: the libfabric back end of NIXL
looks for cxi first, then for efa, and then falls back to TCP \citep{nixl}.

\subsection{Reasons for a Second Interface}

Three reasons led to a second interface beside raw verbs. The first is
development without RDMA hardware. The library has providers, which are back
ends that implement the same calls over different transports, and its TCP and
shared-memory providers run anywhere, so the whole test suite runs on
machines with no RDMA device (\cref{ch:compat}). The second is a
like-for-like baseline: the transport library that the serving engine's
standard connector uses, NIXL \citep{nixl}, has a libfabric back end, and
building on the same library places the two on a common footing. The third is
that published comparisons exist on which to stand. UCX, the other portable
interface, is described by its designers \citep{shamis2015ucx}, its cost over
InfiniBand has been measured per operation \citep{papadopoulou2017ucx}, and
libfabric and UCX have been compared on a production system
\citep{bareford2022} and discussed by an adapter vendor
\citep{smith2023opx}.

\subsection{Objects and Their Verbs Counterparts}

The names in libfabric differ from those in verbs, and most of its objects
have a direct counterpart. \Cref{tab:fimap} gives the mapping
\citep{libfabric_man}.

\begin{table}[htbp]
  \centering
  \caption[libfabric objects and calls beside their verbs counterparts]{libfabric objects and calls beside their verbs counterparts.}
  \label{tab:fimap}
  \small
  \begin{tabular}{L{4.0cm}L{4.6cm}L{5.6cm}}
    \toprule
    Purpose & Verbs & libfabric \\
    \midrule
    Find a device & device list, open device & \code{fi_getinfo}, \code{fi_fabric}, \code{fi_domain} \\
    Group resources & protection domain & domain \\
    Register memory & \code{ibv_reg_mr}, local and remote key & \code{fi_mr_reg}, descriptor and key \\
    Connection endpoint & queue pair & endpoint (\code{fi_endpoint}) \\
    Connection events & a separate connection manager, or keys exchanged by hand & event queue (\code{fi_eq_open}), \code{fi_connect}, \code{fi_listen}, \code{fi_accept} \\
    Report completions & completion queue, \code{ibv_poll_cq} & completion queue, \code{fi_cq_read}, \code{fi_cq_sread} \\
    Count completions & none & counter (\code{fi_cntr_open}) \\
    One-sided write & post send with a write opcode & \code{fi_write}, \code{fi_writedata} \\
    One-sided read & post send with a read opcode & \code{fi_read} \\
    Two-sided message & post send, post receive & \code{fi_send}, \code{fi_inject}, \code{fi_recv} \\
    \bottomrule
  \end{tabular}
\end{table}

Two differences from verbs are worth stating. The first is the type of
endpoint. A message endpoint is connected to one peer, like a reliable
connected queue pair, whereas a reliable-datagram endpoint is not connected to
anyone, and the caller names the destination on each call. On RDMA hardware
libfabric offers the second kind through a utility layer, called rxm, that
manages connected endpoints underneath. The pool supports both, and E4 asks
which is better for each traffic class.

The second difference is the counter, a number that the provider increments
when an operation completes. Reading it costs less than taking an entry off a
completion queue, and it can be configured to count writes that arrive from a
remote peer. It reports how many completed, not which. The next chapter uses this as one of five completion signals.

\section{The RDMA Process Step by Step}
\label{ch:process}

The previous two sections named the objects, and what remains is to see
them act. A single transfer is followed from a cold start to a completed
write and back again, one step at a time, first in libfabric and then in raw verbs, as the reference for
what the pool does on every call.

\subsection{Phase 0: A Link and a Device}

Before any library call there must be an RDMA device on a link that is up,
which means that four conditions hold. The cable is connected and both
interfaces have addresses on the same subnet. An RDMA device exists on that
interface: on the home pair the ConnectX-4 provides it, and on the lab pair
it is created in software (\cref{ch:bringup}). The port of the device is
active and has a RoCEv2 address entry for the interface's IP address, and
this entry, the GID, is what a peer will use to reach it. Finally, the
firewall admits UDP port 4791.

Every failure in this phase shows up later as a connection that never
completes, so the first notebook checks all four before anything else runs.

\subsection{Phase A: Opening the Pool}

Opening the pool builds, for each of the two classes, everything in
\cref{tab:fimap}. The thirteen steps below are the libfabric sequence for one
class on the connecting side. The listening side differs in step 10.

\begin{enumerate}
  \item \code{fi_getinfo} asks what is available. It is called with hints,
        namely the provider, the endpoint type, and the capabilities needed
        (messages, one-sided access, remote completion data), and it returns
        descriptions of matching interfaces.
  \item \code{fi_fabric} opens the network to which the chosen description
        belongs.
  \item \code{fi_domain} opens the device, to which memory registrations and
        endpoints will belong.
  \item \code{fi_eq_open} creates the event queue, where connection events
        will be reported.
  \item \code{fi_cq_open} creates the completion queue for this class, with
        the format that carries remote data, and with a wait object for the
        cold class.
  \item \code{fi_endpoint} creates the endpoint with this class's queue
        depths.
  \item \code{fi_ep_bind} attaches the event queue and the completion queue
        to the endpoint, for transmit and for receive.
  \item \code{fi_enable} makes the endpoint usable.
  \item \code{fi_mr_reg} registers the pool's buffer arena for local use and
        for remote write and read, and returns a descriptor for local use and
        a key for the peer.
  \item The connecting side calls \code{fi_connect} with the peer's address.
        The listening side has opened a passive endpoint, called
        \code{fi_listen}, received a connection request on its event queue,
        created an endpoint from the request and called \code{fi_accept}.
  \item Both sides read their event queue until it reports
        \code{FI_CONNECTED}. Underneath, the provider has walked a queue pair
        through the states of \cref{fig:qpstates}.
  \item Each side sends the other the base address and key of its arena,
        over a small out-of-band TCP connection.
  \item Each side posts receive buffers, because every write with immediate
        data will consume one at the target.
\end{enumerate}

The steps run once per class. After them the pool holds two connected
endpoints, two completion queues and one registered arena that both use.

\subsection{Phase B: One Write}

\Cref{fig:writeseq} follows one write with immediate data from the sender's
call to both completions.

\begin{figure}[htbp]
  \centering
  \begin{tikzpicture}[diagram,
      alt={Sequence diagram with four vertical lifelines, left to right: sender application,
        sender network device, target network device, target application. Message 1, sender
        application to sender device: post a write with immediate data. Message 2, sender device
        to target device: data packets, first, middle and last, the last one carrying the
        immediate value. Message 3, inside the target device: check the key and bounds, place
        the bytes in memory, consume one posted receive. Message 4, target device to sender
        device, dashed: acknowledgment. Message 5, sender device to sender application,
        dashed: completion with status. Message 6, target device to target application, dashed:
        completion with the immediate value. A note under the diagram says the target
        application does nothing until message 6.},
      life/.style={draw=dGray, line width=0.6pt, dashed},
      actor/.style={tint=#1, minimum width=27mm, minimum height=10mm, font=\footnotesize, align=center}]
    \foreach \x/\c/\n/\t in {0/dBlue/a/{Sender\\application}, 3.9/dBlue/b/{Sender\\device}, 7.8/dGreen/c/{Target\\device}, 11.7/dGreen/d/{Target\\application}} {
      \node[actor=\c] (\n) at (\x,0) {\t};
      \draw[life] (\x,-0.5) -- (\x,-7.2);
    }
    \draw[flow] (0,-1.2) -- (3.9,-1.2) node[midway, above, small text] {1. post write + immediate};
    \draw[flow] (3.9,-2.3) -- (7.8,-2.3) node[midway, above, small text] {2. data packets};
    \node[small text, align=center, fill=white] at (7.8,-3.3) {3. check key, place bytes,\\consume one receive};
    \draw[reply] (7.8,-4.4) -- (3.9,-4.4) node[midway, above, small text] {4. acknowledgment};
    \draw[reply] (3.9,-5.5) -- (0,-5.5) node[midway, above, small text] {5. completion, status};
    \draw[reply] (7.8,-6.4) -- (11.7,-6.4) node[midway, above, small text] {6. completion, immediate};
  \end{tikzpicture}
  \caption[One RDMA write with immediate data]{One RDMA write with immediate data. The target application does nothing until step 6.}
  \label{fig:writeseq}
\end{figure}

In detail, the steps are the following, each with the layer that acts.

\begin{enumerate}
  \item The application places the data in a buffer inside the registered
        arena.
  \item The pool chooses the class. For the hot class a payload under the
        provider's inject limit is copied at the call.
  \item The pool posts the write through libfabric with \code{fi_writemsg}
        and the flag for remote completion data, the general form of
        \code{fi_writedata}, and gives it the buffer, its descriptor, the
        32-bit immediate value, the remote address and the remote key.
  \item In the provider, the call becomes a work request on the queue pair's
        send queue. If the queue is full the call returns ``try again'' and
        the pool makes progress on completions before retrying.
  \item On the sender's device, the work request reaches the head of the
        send queue. Work requests on one queue pair are served in order,
        which is fact 1 of \cref{sec:ordering}.
  \item The device cuts the message into packets no larger than the path's
        RDMA packet size, each with a transport header that carries the
        destination queue pair and a sequence number.
  \item On the wire, each packet travels as a UDP datagram inside an IP
        packet inside an Ethernet frame. On Soft-RoCE the kernel builds each
        of these, and on the ConnectX-4 the card does.
  \item For the first packet, the target's device checks that the key is
        valid, that the address range lies inside the registered region, and
        that the region allows remote writes.
  \item The target's device places the payload bytes directly into the
        target's memory. The target application is not scheduled.
  \item The last packet carries the immediate value. The target's device
        takes one posted receive off the receive queue and writes a
        completion that carries the value.
  \item The target's device returns acknowledgments to the sender. A missing
        packet is retransmitted, and after too many retries the queue pair
        enters the error state.
  \item The sender's device writes a completion for the work request. The
        pool reads it with \code{fi_cq_read}, or \code{fi_cq_sread} for the
        cold class, and checks its status.
  \item The target application reads its completion, takes the immediate
        value, and now knows that the bytes are in place (fact 3) and which
        buffer they are in.
\end{enumerate}

Steps 5 to 11 involve no application code on either machine and, on hardware,
no kernel code. That is the whole advantage of RDMA. On Soft-RoCE the same
steps run in the kernel of each machine, which is why the semantics carry over
and the timings do not.

\subsection{Phase C: Read and Send}

An RDMA read runs the same way with the data flowing toward the initiator:
the request goes out, the target's device checks the key and returns the
bytes, and only the initiator receives a completion. By fact 1 a read posted
after a write on the same queue pair is executed after that write, so its
completion proves the write was applied. A send is the two-sided case: the
bytes land in whichever receive buffer the target posted next, and both sides receive a completion. The pool uses send for the hot class's small control
messages and write with immediate data for everything that has a designated
place to land.

\subsection{Phase D: Failure and Teardown}

A queue pair that exhausts its retries, receives a bad key, or loses its peer
enters the error state, and every outstanding work request completes with a
flush status. The pool treats any completion whose status is not success as a
failure of the transfer. It raises an error, records an event in the
telemetry, and never reports a latency for that operation. The connector
above it can fall back to TCP, and records that too (\cref{sec:connector}). Teardown closes objects in the reverse order of Phase A:
endpoints, memory registration, completion and event queues, domain, fabric.

\subsection{The Same Sequence in Raw Verbs}

The verbs implementation performs the same phases with different calls and
with the connection handshake done by hand.

\begin{enumerate}
  \item Open the device and allocate a protection domain.
  \item Register the arena (\code{ibv_reg_mr}).
  \item Create a completion queue and a reliable connected queue pair per
        class.
  \item Pick the address entry on the port that maps to the interface's IPv4
        address.
  \item Exchange queue pair number, starting sequence number, address entry
        and the arena's key with the peer over TCP.
  \item Move the queue pair to initialized, ready-to-receive and
        ready-to-send (\code{ibv_modify_qp}, three times).
  \item Post receives. Post a write with an immediate value
        (\code{ibv_post_send}).
  \item Poll the completion queue (\code{ibv_poll_cq}) and check the status
        of every entry.
\end{enumerate}

Steps 4 to 6 are what libfabric's \code{fi_connect} and its event queue hide,
and they contain a trap on Soft-RoCE: a connection between two machines needs
the address entry that corresponds to a routable address, whereas the default
entry is link-local and works only on a loopback.

\subsection{A Worked Example: One Layer and One Credit}

Take one layer of a 2,048-token prompt, 4 MiB by \cref{eq:kvsize}, moving
from server 1 to server 2 over the direct cable, followed by the 8-byte
credit that server 2 returns. The cold endpoint carries the layer. At the
rate of that link its bytes need 3.36 ms on the wire, and with
4,096-byte RDMA packets it is about a thousand packets (steps 6 to 9, a
thousand times). Server 2's application learns of it once, at step 13.

The credit is an 8-byte write on the hot endpoint, a single packet. If it shared the
cold endpoint and its completion queue, it could be posted behind the next
layer that server 2 is itself sending, and wait milliseconds
(\cref{eq:fifo}). On its own endpoint it is at the head of an empty send
queue. What that is worth in measured microseconds on each pair is reported
in \cref{tab:e3}.

\section{Ultra Ethernet}
\label{sec:uec}

RoCE, the transport of both server pairs, carries InfiniBand's
transport over Ethernet almost unchanged. It expects packets in order and, as
usually deployed, a network that does not drop them, which is obtained with
priority flow control. Flow control of that kind spreads congestion and
blocks traffic behind a paused queue inside the network, and strict ordering
limits the choice of paths \citep{hoefler2025ue}.

The Ultra Ethernet Consortium was formed to replace that transport. A first
group of five companies, AMD, Broadcom, HPE, Intel and Microsoft, began work
in early 2022 \citep{hoefler2025ue}. The consortium was announced on July 19,
2023, as a project of the Joint Development Foundation hosted by the Linux
Foundation, with nine founding members: AMD, Arista, Broadcom, Cisco, Eviden,
HPE, Intel, Meta and Microsoft \citep{uec2023launch}. It had more than 100
member companies by the end of 2024 \citep{hoefler2025ue}. Version 1.0 of its specification was released on June 11, 2025
\citep{uec2025spec10}, and the text consulted here is version 1.0.1, of
September 2025 \citep{uecspec}.

The specification covers the stack from the physical layer to the software
interface. Its center is a new transport, the Ultra Ethernet Transport (UET),
which differs from RoCE in six ways that matter for large clusters
\citep{hoefler2025ue,uecspec,xiong2025ofaue}. It has no connections: a
sender creates the state for a destination with its first data packet, with
no handshake and no queue pair. It offers four kinds of delivery, which are
reliable and ordered; reliable and unordered, the kind meant for bulk data;
reliable and unordered for operations that can safely be applied twice; and
unreliable datagrams. It uses many paths at once, as the sender changes a
field in the header from packet to packet, and switches that hash on that
field spread the packets of one message over many paths. It has two
congestion controls, a sender-based scheme that every device must have and an
optional scheme in which the receiver hands out credits. It defines three
profiles, AI Base, AI Full and HPC, in rising order of features, so that a
device can implement only what its workload needs. Finally, it receives help
from the lower layers, in the form of retry and credit-based flow control on
a single link, optional trimming of packets in switches instead of dropping
them, and optional security.

The part relevant here is the top of that stack. The specification adopts
libfabric as its communication abstraction layer and defines how the calls of
libfabric 2.0 are carried out on UET. It expects each maker of an Ultra
Ethernet endpoint to supply a libfabric provider for it \citep{uecspec}. The
person who chairs the libfabric working group and maintains the library also
takes part in the consortium's software and transport groups
\citep{xiong2025ofaue}. One interface, libfabric, therefore reaches both the
RDMA networks of today, through its verbs provider, and the transport that
the largest network vendors and operators have designed to follow RoCE.

Ultra Ethernet also continues the supercomputing line of \cref{sec:history}:
its specification states that the products which seeded the effort were
inspired by Portals 4, the interface after which libfabric itself was
modeled, and it cites Slingshot for the flow control with which
high-performance networks isolate traffic \citep{uecspec}. Slingshot
separates traffic classes inside the network \citep{desensi2020slingshot},
whereas two servers on one cable have no switch that could do so, which is
why the middleware pursues the same separation at the endpoints.

The next steps that the consortium has announced include
programmable congestion management, congestion signaling carried in packets,
a smaller header, a transport for links inside a system, and collective
operations carried out in switches \citep{uec2024update,uec2025vision}.

\subsection{Requirements for Running the Middleware on Ultra Ethernet}
\label{sec:uec-port}

None of the software described here has run on Ultra Ethernet. No Ultra Ethernet device was
available, and the libfabric tree consulted contains no Ultra Ethernet provider \citep{libfabric_repo}. What can be said, from the
specification and from the pool's source, is what would have to change, and
there are four such points.

The first is the endpoint type. The mapping lists reliable-datagram and
datagram endpoints, and connected endpoints, which the pool uses by default,
do not appear in it. The pool already has a reliable-datagram mode, and E4
compares the two, so E4 is also a first look at the mode Ultra Ethernet would
require. The second is memory registration. The specification requires an
application to register memory with an endpoint, whereas the pool registers
one arena for a whole domain and refuses a provider that asks for
registration per endpoint, which makes this the one change to the code that
is certain. The third is the write with completion data, from which the blob
channel learns of an arrival. The specification requires that call in the AI
Full and HPC profiles and not in AI Base, so the channel would need a device
of one of the two larger profiles, or one of the other signals of
\cref{tab:signals}. The fourth is the library version. The pool asks for
interface version 1.9 and was developed against 1.17, while the mapping is
written against 2.0. The pool builds against release 2.7.0 without a change
to its source, and its tests pass there on the three software providers, so
what remains of this point is the provider itself.

A larger question sits behind these four points. Ultra Ethernet removes ordering and
blocking inside the network, whereas the middleware separates traffic classes
at the endpoints, above the transport, so whether separate endpoints still
matter over a transport that sprays packets across paths is an experiment
for future work. With the interface and its ordering rules in place, the next chapter builds the middleware on them.
